% 分类目录：dl
\documentclass[12pt, a4paper]{article}
\usepackage[margin=2.5cm]{geometry}
\usepackage{ctex}
\usepackage{amsmath, amssymb}
\usepackage{listings}
\usepackage{xcolor}
\usepackage{tabularx}

\lstset{
    language=Python,
    basicstyle=\ttfamily\small,
    keywordstyle=\color{blue},
    commentstyle=\color{green!50!black},
    stringstyle=\color{red},
    showstringspaces=false,
    breaklines=true,
    numbers=left,
    numberstyle=\tiny\color{gray},
    frame=single
}

\title{知识点：扩散模型 (Diffusion Models)}
\author{}
\date{}

\begin{document}
\maketitle

\section{核心定义}
\textbf{中文定义}：扩散模型（Diffusion Models）是一类生成模型，其灵感源自非平衡热力学。它通过定义一个正向扩散过程（逐渐向数据中添加高斯噪声直至其变为纯噪声）和一个反向去噪过程（学习从噪声中恢复原始数据）来学习复杂的数据分布。

\textbf{英文定义}：Diffusion Models are a class of generative models inspired by non-equilibrium thermodynamics. They work by defining a forward diffusion process that systematically adds Gaussian noise to data and a learned reverse process that restores the data from noise.

\textbf{代表作}：DDPM (Denoising Diffusion Probabilistic Models, 2020)。

\section{核心原理：前向与反向过程}
扩散模型的核心由两个互逆的过程组成：

\subsection{前向过程 (Forward Diffusion)}
将原始数据 $x_0$ 逐步转变为纯高斯噪声 $x_T$。每一步添加微小噪声：
$$ q(x_t | x_{t-1}) = \mathcal{N}(x_t; \sqrt{1 - \beta_t} x_{t-1}, \beta_t \mathbf{I}) $$
利用重参数化技巧，可以直接从 $x_0$ 计算任意时刻的 $x_t$：
$$ x_t = \sqrt{\bar{\alpha}_t} x_0 + \sqrt{1 - \bar{\alpha}_t} \epsilon, \quad \epsilon \sim \mathcal{N}(0, \mathbf{I}) $$
其中 $\alpha_t = 1 - \beta_t$，$\bar{\alpha}_t = \prod_{i=1}^t \alpha_i$。

\subsection{反向过程 (Reverse Denoising)}
模型学习预测噪声 $\epsilon_\theta(x_t, t)$，从而实现从 $x_t$ 到 $x_{t-1}$ 的逐步还原：
$$ p_\theta(x_{t-1} | x_t) = \mathcal{N}(x_{t-1}; \mu_\theta(x_t, t), \Sigma_\theta(x_t, t)) $$

\section{训练目标函数}
DDPM 的简化损失函数通常采用均方误差（MSE），目标是让网络预测注入的噪声 $\epsilon$：
$$ L_{simple}(\theta) = \mathbb{E}_{x_0, \epsilon, t} \left[ \| \epsilon - \epsilon_\theta(\sqrt{\bar{\alpha}_t} x_0 + \sqrt{1 - \bar{\alpha}_t} \epsilon, t) \|^2 \right] $$

\section{模型结构：时间敏感型 U-Net}
扩散模型通常使用 U-Net 架构，关键改进是引入了\textbf{时间嵌入 (Time Embedding)}，使模型知道当前处于哪个去噪阶段。
\begin{verbatim}
[ 噪声图像 x_t ] --+--> [ Encoder (下采样) ] --+--> [ Bottleneck ]
                    |           |               |
[ 时间步 t ] -------+--> [ 加入 Time Embedding ] +--> [ Decoder (上采样) ]
                                                |
                                                v
                                         [ 预测的噪声 epsilon ]
\end{verbatim}

\section{关键代码片段 (PyTorch简化实现)}
\begin{lstlisting}
import torch
import torch.nn as nn

class DiffusionTrainer:
    def __init__(self, model, betas):
        self.model = model
        self.alphas = 1.0 - betas
        self.alphas_cumprod = torch.cumprod(self.alphas, dim=0)

    def get_loss(self, x_0):
        batch_size = x_0.shape[0]
        # 1. 随机采样时间步 t
        t = torch.randint(0, len(self.alphas), (batch_size,))
        # 2. 采样标准正态噪声
        noise = torch.randn_like(x_0)
        
        # 3. 计算含噪图像 x_t
        sqrt_alphas_cumprod_t = self.alphas_cumprod[t].sqrt().view(-1, 1, 1, 1)
        sqrt_one_minus_alphas_cumprod_t = (1 - self.alphas_cumprod[t]).sqrt().view(-1, 1, 1, 1)
        x_t = sqrt_alphas_cumprod_t * x_0 + sqrt_one_minus_alphas_cumprod_t * noise
        
        # 4. 模型预测噪声并计算 MSE Loss
        predicted_noise = self.model(x_t, t)
        return torch.nn.functional.mse_loss(noise, predicted_noise)
\end{lstlisting}

\section{深度面试题与解答}
\subsection{基础概念}
\textbf{问：扩散模型相比 GAN 有什么核心优势？}\\
答：1. \textbf{训练极度稳定}：Diffusion 是基于似然概率的显式建模，不存在 GAN 那种难平衡的对抗博弈和模式崩溃。2. \textbf{生成多样性高}：Diffusion 能更好地覆盖数据分布的所有模态。

\subsection{进阶原理}
\textbf{问：什么是 Classifier-Free Guidance (CFG)？为什么它很重要？}\\
答：它是一种在推理时不需要额外分类器就能大幅提升生成质量（文本一致性）的技术。通过同时训练有条件和无条件的模型（即将条件设为空），并在生成时将两者的预测方向进行外推：$\tilde{\epsilon} = \epsilon_{uncond} + w \cdot (\epsilon_{cond} - \epsilon_{uncond})$。

\subsection{工程实战}
\textbf{问：LDM (Latent Diffusion Model) 的贡献是什么？}\\
答：它将扩散过程从像素空间放到了预训练 VAE 的\textbf{潜空间 (Latent Space)}。这极大地降低了计算开销，使得在消费级 GPU 上进行大规模高分辨率图像生成成为可能（Stable Diffusion 的基础）。

\subsection{坑点分析}
\textbf{问：Diffusion 推理速度慢的问题如何解决？}\\
答：原始 DDPM 需要几百上千步迭代。解决方法包括：1. \textbf{DDIM}（确定性采样，仅需 20-50 步）；2. \textbf{蒸馏技术 (Distillation)}，将多步采样能力通过教师模型教给只需几步的弟子模型。

\section{生成模型横向对比表}

\begin{table}[h]
\centering
\begin{tabularx}{\textwidth}{|l|X|X|X|}
\hline
\textbf{维度} & \textbf{VAE} & \textbf{GAN} & \textbf{Diffusion} \\
\hline
训练目标 & 变分下界最大化 & 对抗博弈 & 噪声预测 (MSE) \\
\hline
样本质量 & 模糊 (Over-smoothed) & 高 (可能缺乏多样性) & 极高且多样 \\
\hline
生成速度 & 快 (单次前向) & 快 (单次前向) & 慢 (多次迭代) \\
\hline
稳定性 & 稳定 & 不稳定 & 非常稳定 \\
\hline
\end{tabularx}
\end{table}

\section{参考文献与拓展}
\begin{enumerate}
    \item \textbf{里程碑论文}: Ho, J., et al. (2020). \textit{Denoising Diffusion Probabilistic Models}.
    \item \textbf{Latent Diffusion}: Rombach, R., et al. (2022). \textit{High-Resolution Image Synthesis with Latent Diffusion Models}.
    \item \textbf{进阶资源}: \textit{What are Diffusion Models?} (Lilian Weng's Blog, 深度解析数学原理)。
\end{enumerate}

\end{document}
